\section{Formal kernel of Triadic Modular Holography}
\label{sec:nucleo}

Triadic Modular Holography, abbreviated HMT, begins with a finite structure of positions and operations. The order of construction matters: positions receive arithmetic evaluations, transitions acquire a ternary signature, and the dynamics preserves the information needed to prolong them. Numerical realizations are formed only afterwards. A constant recognized at the end of this process does not replace the rule that produced its region or the history that determines its evaluation.

The name \emph{All Possible Paths}, APP, expresses the role of paths on the support. It is adopted as the name of the construction, with a conceptual reference to the path formulation of quantum mechanics; it does not imply that its arithmetic law is Feynman's path integral. The \emph{Topological Phase Kernel}, TPK,\footnote{The authors also dedicate the initials TPK to Tan Pengkiang. This dedication is independent of the mathematical definition.} denotes the composition of selection, transport, memory update, and prolongation. The acronyms are retained only after the objects they represent have been defined.

\subsection{APP: positional support and arithmetic evaluations}

Consider the nine interior marks on a line graduated from zero to ten,
\[
 D_9=\{1,2,3,4,5,6,7,8,9\}.
\]
Their horizontal and vertical arrangement produces the eighty-one positions of \(D_9^2\). The interior marks must be distinguished from the ten unit intervals of the reference line. This arrangement fixes an alphabet and two directions; it does not yet use a real expansion to determine a trajectory.

Cyclic identification of the positions yields the support
\[
 X_9=(\ZZ/9\ZZ)^2.
\]
The chart of positive representatives retains the symbol nine for the zero class. For every positive integer, define
\begin{equation}
 \rho_9(n)=1+((n-1)\bmod9),\qquad
 q_9(n)=\frac{n-\rho_9(n)}9,
 \qquad n=\rho_9(n)+9q_9(n).
 \label{eq:residuo-cociente}
\end{equation}
The positive digital root is therefore a choice of representative of the residue class. By itself, it does not preserve the integer before reduction; the coordinate \(q_9\) supplies precisely the missing information.

At each position \((i,j)\), accumulating \(i\), repeated \(j\) times, produces \(ij\). The two integer evaluations and their projections are
\[
 S(i,j)=i+j,\qquad P(i,j)=ij,\qquad
 \Sigma(i,j)=\rho_9(i+j),\quad \Pi(i,j)=\rho_9(ij).
\]
The symmetry \((i,j)\mapsto(j,i)\) interchanges the directions of accumulation and preserves both evaluations. This is a property of a single support, not a subsequent identification between two independent matrices.

\begin{definition}[Lifted arithmetic evaluation]
The APP evaluation that preserves both sheets is
\begin{equation}
 \mathcal A(i,j)=
 \bigl((R^+_{ij},q^+_{ij}),(R^\times_{ij},q^\times_{ij})\bigr)
 =\bigl((\rho_9(i+j),q_9(i+j)),(\rho_9(ij),q_9(ij))\bigr).
 \label{eq:app-levantada}
\end{equation}
Here the term \emph{sheet} denotes an evaluation on the same support: additive or multiplicative. The two quotient coordinates belong to the arithmetic state and are not discarded when the residues are reduced.
\end{definition}

\begin{figure}[tbp]
\centering
\begin{tikzpicture}[x=.53cm,y=.53cm,font=\small]
\fill[black!5] (1.5,-.5) rectangle (2.5,8.5);
\fill[black!5] (4.5,-.5) rectangle (5.5,8.5);
\fill[black!5] (-.5,2.5) rectangle (8.5,3.5);
\fill[black!5] (-.5,5.5) rectangle (8.5,6.5);
\fill[black!12] (7.5,-.5) rectangle (8.5,8.5);
\fill[black!12] (-.5,-.5) rectangle (8.5,.5);
\foreach \i in {1,...,9}{
 \node at (\i-1,9.05) {\(\i\)};
 \node at (-1.05,9-\i) {\(\i\)};
 \foreach \j in {1,...,9}{
  \pgfmathtruncatemacro{\v}{mod(\i*\j-1,9)+1}
  \node at (\j-1,9-\i) {\(\v\)};
 }
}
\foreach \k in {0,...,9}{
 \draw[black!25] (\k-.5,-.5)--(\k-.5,8.5);
 \draw[black!25] (-.5,\k-.5)--(8.5,\k-.5);
}
\draw[line width=.65pt] (2.5,3.5) rectangle (4.5,5.5);
\node[anchor=west,align=left,text width=6.5cm] at (10,7.2) {One support: \(81\) positions.\\[5pt]Two evaluations:\\ \(S(i,j)=i+j\), \(P(i,j)=ij\).\\[5pt]Each row is generated by accumulation:\\ \(P(i,j+1)-P(i,j)=i\).};
\node[anchor=west,align=left,text width=6.5cm] at (10,1.8) {The chart represents \(0\pmod9\) by \(9\).\\[5pt]The highlighted block is\\ \(B_c=\left(\begin{smallmatrix}7&2\\2&7\end{smallmatrix}\right)\),\\ with eigenvalues \(9\) and \(5\).};
\end{tikzpicture}
\caption{APP multiplication table in positive representatives. Non-unit rows and columns are distinguished by gray shading; the frame identifies the adjacent diagonal block with equal diagonals. The displayed table is an observable of the arithmetic torus; its representation edges do not constitute an intrinsic topological boundary of the torus.}
\label{fig:app}
\end{figure}


\begin{proposition}[Reconstruction of the evaluations]
The two pairs in \eqref{eq:app-levantada} determine \(i+j\) and \(ij\) exactly. The residues alone do not determine these integer evaluations or their quotients.
\end{proposition}
\begin{proof}
The first statement is \eqref{eq:residuo-cociente} applied to each sheet. For the second, \(10=1+9\) and \(19=1+2\cdot9\) have the same residue and different quotients. On the two-dimensional support, positions \((5,5)\) and \((8,2)\) give the same sum and the same product residue, but products \(25=7+18\) and \(16=7+9\). The residue projection identifies data that the lifted evaluation distinguishes.
\end{proof}

\subsubsection{Cayley graph, toroidal topology, and winding number}

The cardinal translations are given, in row and column coordinates, by
\[
 \nu_N=(-1,0),\quad\nu_E=(0,1),\quad
 \nu_S=(1,0),\quad\nu_O=(0,-1),\qquad T_d(x)=x+\nu_d.
\]
The first coordinate increases southwards and the second eastwards. The graph
\[
 \mathscr G_9=\operatorname{Cay}(X_9,\{\nu_N,\nu_E,\nu_S,\nu_O\})
\]
is the two-dimensional cyclic lattice: it has eighty-one vertices and four outgoing oriented edges per vertex. It is a Cayley graph for the \emph{additive} structure of the support. Multiplication of residues has zero divisors and does not turn all APP rows into a multiplicative group.

A cardinal word \(w=d_1\cdots d_r\) and an origin \(x_0\) determine
\[
 x_k=x_0+\sum_{j=1}^k\nu_{d_j}\pmod9.
\]
There are \(81\cdot4^r\) oriented paths of length \(r\), including returns and backtracking. When the path closes, its integer displacement before reduction defines
\begin{equation}
 \operatorname{wind}(w)=\frac19
 \bigl(\#S(w)-\#N(w),\#E(w)-\#O(w)\bigr)\in\ZZ^2.
\end{equation}
Path reversal changes the sign of this vector. Returning to the same residue position therefore does not require the winding to vanish. This is the first example of a distinction that will reappear in phase return: one coordinate may close while another records a nontrivial evolution.

\subsubsection{Transport quotient and cocycle identity}

Let \(s_9:\ZZ/9\ZZ\to D_9\) be the section of positive representatives. The transport quotient associated with additive composition, also called the \emph{carry} in positional arithmetic, is
\[
 c_9(a,b)=\frac{s_9(a)+s_9(b)-s_9(a+b)}9.
\]
The compositional character of this memory is expressed by an exact identity.

\begin{lemma}[Cocycle identity]
For \(a,b,c\in\ZZ/9\ZZ\),
\[
 c_9(a,b)+c_9(a+b,c)=c_9(b,c)+c_9(a,b+c).
\]
\end{lemma}
\begin{proof}
Both sums equal
\[
 \frac{s_9(a)+s_9(b)+s_9(c)-s_9(a+b+c)}9.
\]
Associativity of the sum of lifted representatives requires precisely this additional term when composing in the quotient.
\end{proof}

With positive representatives, this cocycle is not normalized at the identity: \(c_9(0,a)=1\). If \(s_0\) takes representatives in \(\{0,\ldots,8\}\), its normalized carry \(c_0\) satisfies
\[
 c_9(a,b)=c_0(a,b)+\mathbf1_{a=0}+\mathbf1_{b=0}-\mathbf1_{a+b=0}.
\]
The difference is the change of section; it does not alter the cocycle identity. The subsequent decomposition of the calendar into phase and window will use representatives from zero to eight.

The totals of both evaluations distinguish the visible contribution from that retained in the quotient:
\[
 \sum S=810,\quad\sum P=2025,\quad
 \sum\Sigma=405,\quad\sum\Pi=459,
\]
\[
 \sum q^+=45,\qquad\sum q^\times=174.
\]
It follows that
\begin{equation}
 2025-810=(459-405)+9(174-45)=54+9\cdot129.
 \label{eq:defecto-integrado}
\end{equation}
The sum--product difference is not exhausted by the visible defect \(54\). The identity also preserves the quotient defect \(129\). Removing it would change the arithmetic being transported, even if the residue table remained identical.

% BEGIN R27_MG03
\subsubsection{Modulus refinement and preservation of the two readers}
\label{mg03:refinamiento}

Identity \eqref{eq:defecto-integrado} preserves the integer defect
\(1215\) through residue \(54\) and quotient \(129\). Its
generalization requires separate specifications of the reading modulus,
the support size and the number of coordinates. These three operations
produce distinct counting laws from the same APP sheets.

For \(N\geq2\), fix the positive representative
\[
 \rho_N(s)=1+((s-1)\bmod N)\in\{1,\ldots,N\},
 \qquad s=\rho_N(s)+NQ_N(s),\quad Q_N(s)\in\mathbb Z.
\]
On \(\{1,\ldots,N\}^2\), the full-modulus readings are
\(\rho_N(i+j)\) and \(\rho_N(ij)\). Their sums are denoted by
\(S_+^{\mathrm{full}}(N)\) and \(S_\times^{\mathrm{full}}(N)\).

\begin{proposition}[Defect law for an integer modulus]
For every \(N\geq2\),
\begin{align}
 S_+^{\mathrm{full}}(N)&=\frac{N^2(N+1)}2,
 \label{mg03:suma-plena}\\
 S_\times^{\mathrm{full}}(N)&=
 \frac N2\left(N^2+\sum_{i=1}^{N}\gcd(i,N)\right),
 \label{mg03:producto-pleno}\\
 \Delta_N^{\mathrm{full}}&=
 \frac N2\left(\sum_{i=1}^{N}\gcd(i,N)-N\right).
 \label{mg03:defecto-pleno}
\end{align}
\end{proposition}

\begin{proof}
In an additive row, translation by \(i\) permutes the \(N\) residue
classes and their positive representatives. Each row sums to
\(N(N+1)/2\), proving \eqref{mg03:suma-plena}. For a multiplicative
row, let \(g=\gcd(i,N)\). The image of \(j\mapsto ij\pmod N\)
contains the \(N/g\) residue classes divisible by \(g\), and each
has \(g\) preimages. Their positive representatives are
\(g,2g,\ldots,N\), whose sum is
\[
 g\frac{(N/g)(N/g+1)}2=\frac{N(N+g)}{2g}.
\]
Multiplicity \(g\) gives \(N(N+g)/2\) per row. Summation over
\(i\) yields \eqref{mg03:producto-pleno}; subtracting
\eqref{mg03:suma-plena} gives \eqref{mg03:defecto-pleno}.
\end{proof}

\begin{corollary}[Prime-power moduli]
If \(N=p^m\), where \(p\) is prime and \(m\geq1\), then
\begin{equation}
 \sum_{i=1}^{p^m}\gcd(i,p^m)
 =p^m+m(p-1)p^{m-1},\qquad
 \Delta_{p^m}^{\mathrm{full}}
 =\frac{m(p-1)}2p^{2m-1}.
 \label{mg03:primo-potencia}
\end{equation}
\end{corollary}

\begin{proof}
For each \(0\leq k<m\), there are
\(p^{m-k}-p^{m-k-1}\) indices with exact \(p\)-adic valuation
\(k\). Their greatest common divisor with \(p^m\) is \(p^k\),
so this stratum contributes \((p-1)p^{m-1}\). The \(m\) strata
contribute \(m(p-1)p^{m-1}\), whereas the index \(i=p^m\)
contributes \(p^m\). Substitution into \eqref{mg03:defecto-pleno}
gives the second formula.
\end{proof}

The triadic specialization is
\(\Delta_{3^m}^{\mathrm{full}}=m3^{2m-1}\). The first moduli
display both sums together with their difference:
\[
\begin{array}{c|r|r|r}
 N&S_+^{\mathrm{full}}&S_\times^{\mathrm{full}}&\Delta_N^{\mathrm{full}}\\\hline
 3&18&21&3\\
 9&405&459&54\\
 27&10206&10935&729\\
 81&269001&277749&8748
\end{array}
\]

Periodic refinement instead retains \(\rho_9\) and enlarges the
support to \(\{1,\ldots,9r\}^2\), with integer \(r\geq1\). Each pair of residues modulo
nine has \(r^2\) representatives in this support, for both addition
and multiplication. Consequently,
\begin{equation}
 S_+^{\mathrm{per}}(9r;9)=r^2\,405,\quad
 S_\times^{\mathrm{per}}(9r;9)=r^2\,459,\quad
 \Delta^{\mathrm{per}}(9r;9)=r^2\,54.
 \label{mg03:periodico}
\end{equation}
For \(r=3\), the sums are \(3645\) and \(4131\), and the defect
is \(486\). The full-modulus reading at \(27\) produces \(729\)
because it updates both the modulus and its representatives; the periodic
reading produces \(486\) because it retains modulus nine and multiplies
its incidences by nine. In both readings, the divisions
\(i+j=R^++NQ^+\), \(ij=R^\times+NQ^\times\), with the modulus
actually employed, preserve the integer identity at each cell and after
summation. This specifies the datum accompanying each defect.

\subsubsection{Dimensional variation at fixed modulus nine}
\label{mg03:dimension}

A second extension retains the modulus and increases the number of
factors. For \(d\geq1\), let
\[
 \mathcal A_d=\{1,\ldots,9\}^d,\qquad
 \Sigma_d(x)=\rho_9\!\left(\sum_{j=1}^d x_j\right),\qquad
 \Pi_d(x)=\rho_9\!\left(\prod_{j=1}^d x_j\right).
\]

\begin{proposition}[Dimensional count for accumulation and multiplication]
The sums of the two readers over \(\mathcal A_d\) are
\begin{equation}
 S_{+,d}=45\,9^{d-1},\qquad
 S_{\times,d}=9^{d+1}-9(d+3)6^{d-1}.
 \label{mg03:ley-dimensional}
\end{equation}
Therefore,
\begin{equation}
 \frac{S_{\times,d}-S_{+,d}}{9^d}
 =4-(d+3)\left(\frac23\right)^{d-1}
 \longrightarrow4.
 \label{mg03:limite-dimensional}
\end{equation}
\end{proposition}

\begin{proof}
For each fixed choice of the first \(d-1\) coordinates, the final
coordinate traverses every additive residue once. Each positive
representative occurs \(9^{d-1}\) times, and their sum is \(45\);
this proves the additive formula.

For multiplication, distinguish the units \(U=\{1,2,4,5,7,8\}\),
the multiples of three of valuation one \(T=\{3,6\}\), and the
zero representative \(9\). The \(6^d\) words whose factors all
belong to \(U\) distribute their products uniformly over the six
units: fixing \(d-1\) factors leaves a permutation of the group
under multiplication by the last factor. Their contribution is
\(27\,6^{d-1}\).

The words with exactly one factor in \(T\) and all other factors
in \(U\) number \(2d6^{d-1}\). For each position of the singular
factor and each choice of the remaining units, its two possible values
\(3,6\) produce each of the residues \(3,6\) once. Their contribution
is \(9d6^{d-1}\). Every remaining word contains a factor \(9\),
or at least two factors in \(T\), so its product has representative
\(9\). Their number is \(9^d-6^d-2d6^{d-1}\). Thus
\[
 S_{\times,d}=27\,6^{d-1}+9d6^{d-1}
 +9\bigl(9^d-6^d-2d6^{d-1}\bigr),
\]
which simplifies to \eqref{mg03:ley-dimensional}. Dividing by
\(9^d\) and subtracting the additive mean \(5\) gives
\eqref{mg03:limite-dimensional}; the term
\((d+3)(2/3)^{d-1}\) tends to zero.
\end{proof}

Dimension two recovers \(S_{+,2}=405\), \(S_{\times,2}=459\)
and total defect \(54\). The limit \(4\) belongs to the mean
defect under increasing dimension. Total defect, preserved quotient
and dimensional mean are specified readers of APP arithmetic; TRIT
and TPK subsequently transport these data with their orientation,
modulus and mathematical residence.

% END R27_MG03

\subsubsection{Ternary reduction and multiplicative radical}

The ring \(R=\ZZ/9\ZZ\) has unit group \(\{1,2,4,5,7,8\}\) and maximal ideal
\[
 \mathfrak m=(3)=\{0,3,6\},\qquad\mathfrak m^2=0.
\]
Every additive row is a permutation of the nine representatives. A multiplicative row is such a permutation exactly when its index is a unit. Rows three and six contain three repeated values; row nine contains only the representative of the zero class. Thus multiplicative folding already distinguishes the units from the nilpotent sector before any continuous geometry is introduced.

The map \(x\mapsto3x\) has kernel and image both equal to \(\mathfrak m\). It induces an isomorphism of vector spaces over \(\FF_3\),
\[
 R/\mathfrak m\longrightarrow\mathfrak m,\qquad [x]\longmapsto3x,
\]
but not a ring isomorphism. In particular, the visible sequence \(3,6,9\) admits two different readings: its direct reduction modulo three is \(0,0,0\), whereas extracting the factor three and then reducing its internal coordinate gives \(1,2,0\). This second operation preserves a coordinate of the ideal; it is not the same projection as directly reducing the original number.

\subsubsection{Latin realization of the additive evaluation}

The difference between the two evaluations can also be expressed through
operators. Let \((e_r)_{r\in\ZZ/9\ZZ}\) be the orthonormal basis of \(\CC^9\).
The assignment \(v_{ab}=e_{a+b}\) produces an orthonormal basis in every row
and column. The projectors \(P_{ab}=v_{ab}v_{ab}^{*}\) satisfy
\[
 \sum_bP_{ab}=I_9=\sum_aP_{ab},\qquad P_{ab}^2=P_{ab}.
\]
Indeed, each translation \(b\mapsto a+b\) permutes the nine residues.
This yields a realization of a Latin square by rank-one projectors,
in the terminology of quantum Latin squares
\cite{mustovicary,delascuevas2023}. This construction starts from the
additive sheet already defined and makes explicit the map to its operator representation.

The multiplicative evaluation preserves different content: for \(a=3\), the
vectors \(e_{3b}\) repeat the residues \(0,3,6\), and
\[
 \sum_{b\in\ZZ/9\ZZ}|e_{3b}\rangle\langle e_{3b}|
 =3\bigl(|e_0\rangle\langle e_0|+|e_3\rangle\langle e_3|
                    +|e_6\rangle\langle e_6|\bigr).
\]
The multiplicity is recorded in the operator and characterizes the multiplicative
radical. Thus the toroidal support, the Latin property, and the operator
resolution of the identity are properties requiring different definitions.
The quantum magic squares studied by De las Cuevas, Drescher, and Netzer
require positive entries and row and column sums equal to the identity;
their theory also distinguishes dilations with commuting entries~\cite{delascuevas2020}.
This reference provides a precise language for comparing realizations of APP.

\subsection{TRIT: orientation and quadratic regimes}

\begin{definition}[Ternary signature and local lift]
The TRIT signature uses the balanced identification
\[
 \FF_3\longrightarrow\TRIT,\qquad 0\mapsto0,\quad1\mapsto+1,\quad2\mapsto-1.
\]
For \(n\in\ZZ\), its lift preserves \(n=r+3c\), where \(r\in\TRIT\) and \(c\in\ZZ\). In a transition of bounded local amplitude, the three states \((0,-1),(0,0),(0,+1)\) may exhibit the same zero residue and different carries.
\end{definition}

The signature does not replace the trajectory. A zero residue does not mean an absence of orientation, phase, or quotient information. Nor does it introduce a physical quantity by itself. Its function is to distinguish algebraically the transport classes that will be realized next.

\subsubsection{Three quadratic regimes and a common exponential}

Once the ternary signature has been obtained, its quadratic realization is defined by
\begin{equation}
 \mathbb A_\tau=\RR[X]/(X^2+\tau),\qquad \tau\in\TRIT.
 \label{eq:algebras-trit}
\end{equation}
If \(J_\tau\) is the class of \(X\), then \(J_\tau^2=-\tau I\). The structure distinguishes the elliptic type \(\tau=+1\), the parabolic type \(\tau=0\), and the hyperbolic type \(\tau=-1\). The real realization does not generate the signature: it expresses the types already produced by oriented reduction.

\begin{proposition}[Exponentiation of the three types]
In each algebra \eqref{eq:algebras-trit},
\begin{equation}
 \exp(tJ_\tau)=C_\tau(t)I+S_\tau(t)J_\tau,
\end{equation}
where
\[
 C_\tau(t)=\sum_{n\ge0}\frac{(-\tau)^nt^{2n}}{(2n)!},\qquad
 S_\tau(t)=\sum_{n\ge0}\frac{(-\tau)^nt^{2n+1}}{(2n+1)!}.
\]
The three realizations are
\[
 \begin{array}{c|c|c}
 \tau&J_\tau^2&\exp(tJ_\tau)\\\hline
 +1&-I&\cos t\,I+\sin t\,J_\tau\\
 0&0&I+tJ_\tau\\
 -1&I&\cosh t\,I+\sinh t\,J_\tau.
 \end{array}
\]
\end{proposition}
\begin{proof}
Separating the even and odd terms of the exponential series and substituting \(J_\tau^{2n}=(-\tau)^nI\) and \(J_\tau^{2n+1}=(-\tau)^nJ_\tau\) gives the formulas. For \(\tau=0\), all terms of degree at least two vanish. In the other two cases, the trigonometric or hyperbolic series are recovered.
\end{proof}

The exponential is common to all three regimes: it is not assigned exclusively to the parabolic branch. The invariants
\[
 C_\tau(t)^2+\tau S_\tau(t)^2=1
\]
follow from \(\exp(tJ_\tau)\exp(-tJ_\tau)=I\). The same algebraic principle preserves unity through realizations that differ in their quadratic type.

In the hyperbolic regime, \(P_\pm=(I\pm J_{-1})/2\) are complementary projectors and
\[
 \exp(tJ_{-1})=e^tP_++e^{-t}P_-.
\]
Growth and contraction appear jointly in two eigendirections. In the elliptic regime, the parameter describes a rotation; in the parabolic regime, a nilpotent translation. The temporal interpretation adopted by HMT associates \(+1\) with return with memory, \(0\) with the threshold, and \(-1\) with bidirectional opening. This interpretation requires an ordering of the parametrization; it does not add a fourth algebraic type.

% Cumulative TRIT extension; inserted before the regime figure.
% It neither replaces nor removes the preceding development.
\subsubsection{Balanced division and composition of lifts}
\label{trit-ampl-div-balanceada}

Ternary reduction expresses an integer evaluation of APP through a visible
coordinate and a prolongation coordinate. Its function is not merely to
abbreviate the notation for an integer: it provides a composition rule
determining how the quotient changes when evaluations are added or multiplied.
The balanced choice also makes the compatibility of this rule with sign
reversal explicit.

\begin{definition}[Balanced ternary coordinates]
\label{trit-ampl-def-coordenadas}
For \(n\in\mathbb Z\), let
\[
 q_3(n)=\left\lfloor\frac{n+1}{3}\right\rfloor,\qquad
 r_3(n)=n-3q_3(n).
\]
The map
\[
 \mathcal B:\mathbb Z\longrightarrow\TRIT\times\mathbb Z,\qquad
 n\longmapsto(r_3(n),q_3(n))
\]
is called the balanced ternary decomposition. Its inverse is
\(\mathcal L(r,c)=r+3c\).
\end{definition}

Indeed, \(r_3(n)\in\{-1,0,+1\}\). If two pairs represent the same integer,
\(r+3c=s+3d\), then \(r-s\) is a multiple of \(3\) and belongs to the interval
\([-2,2]\); necessarily \(r=s\) and \(c=d\). The representation is therefore
unique. Division can be repeated on \(q_3(n)\), so that each level preserves
the data needed to reconstruct the preceding one. The quotient is not a
perturbation of the residue: it is its complementary coordinate.

The signature obtained from an evaluation \(F\), with \(F=S\) or \(F=P\) in APP,
is \(r_3\circ F\). When reduction acts on the phase cursor, the same balanced
section produces the classes \(\{1,4,7\}\), \(\{2,5,8\}\), and
\(\{3,6,9\}\), with values \(+1,-1,0\). The TPK selector defined below
uses these classes to determine the operating mode.
Reading an evaluation and selecting a mode thus share the same ternary
reduction, but have different domains and functions.

\Needspace{16\baselineskip}
\begin{proposition}[Arithmetic of balanced pairs]
\label{trit-ampl-prop-aritmetica}
Let \(n=r+3c\) and \(m=s+3d\), with \(r,s\in\TRIT\). Define
\[
 \chi(r,s)=\frac{r+s-r_3(r+s)}3.
\]
Integer operations are transported to the pairs by
\begin{align}
 (r,c)\boxplus(s,d)
 &=\bigl(r_3(r+s),c+d+\chi(r,s)\bigr),
 \label{trit-ampl-eq-suma}\\
 (r,c)\boxtimes(s,d)
 &=\bigl(rs,rd+sc+3cd\bigr).
 \label{trit-ampl-eq-producto}
\end{align}
The map \(\mathcal L\) preserves both operations.
\end{proposition}
\begin{proof}
The equality \(r+s=r_3(r+s)+3\chi(r,s)\) gives
\[
 n+m=r_3(r+s)+3\bigl(c+d+\chi(r,s)\bigr).
\]
On the other hand,
\[
 nm=rs+3(rd+sc+3cd).
\]
Since the product of two balanced representatives again belongs to
\(\{-1,0,+1\}\), it requires no additional residue reduction. Uniqueness
of the decomposition proves both formulas.
\end{proof}

For example, \(2=(-1)+3\) and \(4=1+3\). Their sum and product are reconstructed
as
\[
 (-1,1)\boxplus(1,1)=(0,2),\qquad
 (-1,1)\boxtimes(1,1)=(-1,3).
\]
The resulting evaluations are \(6\) and \(8\). Adding the quotients suffices
in this additive example; the product involves the cross terms
\(rd\), \(sc\), and \(3cd\). The discrepancy between the two operations
also resides in the prolongation coordinate, not only in the residue.

The term \(\chi\) is a cocycle of the balanced section. If
\(r\oplus s=r_3(r+s)\), it satisfies
\[
 \chi(r,s)+\chi(r\oplus s,t)
 =\chi(s,t)+\chi(r,s\oplus t).
\]
Both sides represent the quotient of the same sum of three integers.
This identity ensures that regrouping the operations does not change the
reconstructed evaluation. At the same time, projection onto the first
component suppresses this accounting: \(1+1=-1+3\), whereas the residue
reading retains only \(-1\). The information removed is determined,
in this case, by \(\chi(1,1)=1\).

The relation with modulus \(9\) preserves a second discrete scale.
The map
\[
 \beta:\TRIT\times\mathbb F_3\longrightarrow\mathbb Z/9\mathbb Z,
 \qquad (r,\bar c)\longmapsto r+3c\pmod9
\]
is well defined and bijective. Changing the representative of \(\bar c\)
adds a multiple of \(9\); conversely, equality of images fixes
\(r\) first, by reduction modulo \(3\), and then \(\bar c\). Addition
transported by this bijection contains the term
\(\chi(r,s)\) from \eqref{trit-ampl-eq-suma}, now reduced modulo \(3\).
Thus the two ternary levels do not compose by independent coordinatewise
addition: normalization of the first level modifies the second.

In particular, the classes \(3\), \(6\), and \(0\) of
\(\mathbb Z/9\mathbb Z\) have first component \(0\), but second components
\(1\), \(2\), and \(0\), respectively. The distinction already observed
between direct reduction and extraction of the ideal coordinate appears
as a property of \(\beta\). Preserving this coordinate prevents the zero
residue sector from being artificially collapsed into a single position.
Modulus \(1000\), used later to organize decimal blocks,
serves a different positional-representation function; its quotient is
preserved through the corresponding division. A change of base preserves
the integer when it retains residue and quotient, whereas coincidence of
residues in different bases does not by itself identify their transport states.

\Needspace{11\baselineskip}
\subsubsection{Inversion, residue class, and non-visible coordinates}
\label{trit-ampl-inversion}

Additive inversion lifts unambiguously:
\[
 \mathcal B(-n)=(-r_3(n),-q_3(n)).
\]
In particular, \(3\), \(0\), and \(-3\) have coordinates
\((0,1)\), \((0,0)\), and \((0,-1)\). They share a zero signature but
represent different evaluations. In a transition restricted to a quotient
of maximum amplitude \(1\), these are the three possibilities of the
visible fiber over zero; in a general integer lift, that fiber contains
all pairs \((0,c)\), with \(c\in\mathbb Z\).

The quotient coordinate recovers the evaluated integer, but does not by
itself recover the trajectory that produced it. Two APP paths may share
evaluation, residue, and quotient while differing in orientation, sequence
of modes, or winding. Thus decomposition \(\mathcal B\) is integrated into
the transport state rather than replacing it. The arithmetic memory of a
division and the memory of a trajectory are related coordinates, not synonyms.

For a signature sequence of length \(n\), oriented reversal is
\[
 C_{\rm rev}(\tau_1,\ldots,\tau_n)
   =(-\tau_n,\ldots,-\tau_1).
\]
Applying it twice restores the sequence. The operation reverses the order
of positions and the sign of each signature; it preserves the length and
the zero positions in reversed order. In the adopted parametrization,
it interchanges return with memory, associated with \(+1\), and bidirectional
opening, associated with \(-1\), while retaining the threshold \(0\).
Transport of the remaining trajectory data is specified by the corresponding
TPK involution.

This signature reversal must be distinguished from conjugation within a
quadratic algebra. The former may change the type \(\tau\); the latter,
constructed below, acts within a type already fixed. This distinction
prevents attributing an algebraic equivalence between different regimes
to a sign reversal.

\subsubsection{Quadratic product, conjugation, and algebraic norm}
\label{trit-ampl-norma}

An element of \(\mathbb A_\tau\) has a unique expression \(aI+bJ_\tau\).
The relation \(J_\tau^2=-\tau I\) determines its product:
\begin{equation}
 (aI+bJ_\tau)(cI+dJ_\tau)
   =(ac-\tau bd)I+(ad+bc)J_\tau.
 \label{trit-ampl-producto-cuadratico}
\end{equation}
This is a real realization of the type already selected. The coordinates
\(a,b\) are not APP residues, and the product
\eqref{trit-ampl-producto-cuadratico} does not replace
\eqref{trit-ampl-eq-producto}; it expresses another operation in the
declared quadratic codomain.

\begin{definition}[Conjugation and algebraic norm]
\label{trit-ampl-def-norma}
Conjugation of type \(\tau\) and its algebraic norm are
\[
 \overline{aI+bJ_\tau}=aI-bJ_\tau,\qquad
 N_\tau(aI+bJ_\tau)=a^2+\tau b^2.
\]
\end{definition}

\begin{proposition}[Multiplicativity and invertibility]
\label{trit-ampl-prop-norma}
For \(z,w\in\mathbb A_\tau\),
\[
 z\bar z=N_\tau(z)I,\qquad
 N_\tau(zw)=N_\tau(z)N_\tau(w).
\]
Moreover, \(z\) is invertible if and only if \(N_\tau(z)\ne0\), in which case
\[
 z^{-1}=\frac{\bar z}{N_\tau(z)}.
\]
\end{proposition}
\begin{proof}
The product of \(aI+bJ_\tau\) and \(aI-bJ_\tau\) is
\((a^2+\tau b^2)I\). Conjugation is multiplicative and the algebra is
commutative, so that
\((zw)\overline{zw}=(z\bar z)(w\bar w)\). If the norm does not vanish,
the stated formula supplies the inverse. Conversely, if \(zw=I\),
multiplicativity implies \(N_\tau(z)N_\tau(w)=1\).
\end{proof}

The quadratic form \(N_\tau\) is positive definite only in the elliptic type.
For \(\tau=0\), it equals \(a^2\), and the element \(J_0\) is nonzero although
its square and norm are zero. For \(\tau=-1\), it equals \(a^2-b^2\);
the nonzero elements \(I+J_{-1}\) and \(I-J_{-1}\) have zero norm and
their product vanishes. Here the algebraic norm is a multiplicative
quadratic form: it does not presuppose a positive vector-space norm.

The three types admit faithful realizations by
\[
 aI+bJ_\tau\longmapsto
 \begin{pmatrix}a&-\tau b\\b&a\end{pmatrix}.
\]
The determinant of this matrix is \(N_\tau(aI+bJ_\tau)\).
The complex structure, the dual-number structure, and the split real
structure are thus represented by a single matrix expression.
Each preserves its own law and zero divisors. The form of signature
\((1,1)\) in the last case admits a Lorentzian interpretation in dimension
\(1+1\); identification of a physical metric additionally requires a
specific space, field, and transport rule.

\subsubsection{Multiplicative conservation and composition of evolutions}
\label{trit-ampl-evoluciones}

The exponential already obtained belongs to the set of norm-one elements:
\[
 N_\tau\bigl(\exp(tJ_\tau)\bigr)=1.
\]
This conservation persists under composition of parameters. Indeed, all
factors are functions of the same generator, and
\[
 \exp(tJ_\tau)\exp(sJ_\tau)=\exp((t+s)J_\tau).
\]
Comparing the two components gives the addition laws
\begin{align}
 C_\tau(t+s)&=C_\tau(t)C_\tau(s)-\tau S_\tau(t)S_\tau(s),\\
 S_\tau(t+s)&=S_\tau(t)C_\tau(s)+C_\tau(t)S_\tau(s).
\end{align}
Conjugation changes \(t\) to \(-t\) and coincides with inversion within
this subgroup. It does not change \(\tau\): it reverses an evolution of
the chosen regime.

In the elliptic type, the preserved form is \(a^2+b^2\), and the exponential
trajectory traverses the unit-norm circle. In the hyperbolic type, the
eigencoordinates are \(u=a+b\), \(v=a-b\), with \(uv=N_{-1}(z)\).
The exponential gives
\[
 (u,v)=(e^t,e^{-t}),\qquad uv=1.
\]
Growth of one coordinate is accompanied by contraction of the other.
Conservation expresses a relation between the two coordinates; it does
not assert that each remains constant.

In the parabolic type,
\[
 (I+tJ_0)(I+sJ_0)=I+(t+s)J_0,\qquad
 (I+tJ_0)^{-1}=I-tJ_0.
\]
The evolution does not stop because \(J_0^2=0\). Nilpotency removes terms
of order at least two, but preserves the linear term that transports
the parameter. The visible symbol \(0\), the zero element of the algebra,
and the nonzero generator \(J_0\) are distinct. This distinction is
necessary to describe the threshold without identifying it with an absence
of state or dynamics.

\begin{proposition}[Composition of affine coordinates]
\label{trit-ampl-prop-pendientes}
In the chart where the scalar component does not vanish, represent a
direction by \(I+uJ_\tau\). If \(1-\tau uv\ne0\), the direction of its
product with \(I+vJ_\tau\) has coordinate
\[
 u\star_\tau v=\frac{u+v}{1-\tau uv}.
\]
\end{proposition}
\begin{proof}
We have
\[
 (I+uJ_\tau)(I+vJ_\tau)
   =(1-\tau uv)I+(u+v)J_\tau.
\]
Dividing by the scalar component gives the stated expression.
\end{proof}

The affine chart preserves the ratio of components, not the scalar factor
divided out. When \(1-\tau uv=0\), the product must be examined before
changing charts: it may have a well-defined direction with zero scalar
component, or be the zero element in the presence of zero divisors.
The composition law still belongs to the full algebra; a coordinate
singularity does not automatically imply a singularity of the object.
This formula will reappear in the elliptic regional reading. Its use
requires retaining the denominator and, where applicable, the normalization factor.

\subsubsection{Transport of type and subsequent realizations}
\label{trit-ampl-puentes}

The algebras \(\mathbb A_{+1}\), \(\mathbb A_0\), and
\(\mathbb A_{-1}\) are not isomorphic as unital real algebras.
The first is a field; the second contains a nonzero nilpotent; the
third has nontrivial idempotents and is isomorphic to
\(\mathbb R\times\mathbb R\). Thus a change of signature organizing a
route cannot be interpreted as an undifferentiated isomorphism between
the three realizations. The TPK preserves the regime index, the domains,
and the effective operator acting at each transition.

\subsubsection{Commutator norm of the arithmetic projections}

The electronic prefactor is obtained without using \(\alpha\), torsion,
or a dimensional unit. The domain is the third realization of APP,
\(\mathcal A_3=\{1,\ldots,9\}^3\), equipped with the uniform measure.
The function space is \(\ell^2(\mathcal A_3;\mathbb R)\), with inner
product
\[
 \langle f,g\rangle=\frac1{729}\sum_{x\in\mathcal A_3}f(x)g(x).
\]
In this block, \([P,Q]=PQ-QP\) denotes the additive operator commutator.
We define
\[
 \Sigma(x)=\rho_9(x_1+x_2+x_3),\qquad
 \Pi(x)=\rho_9(x_1x_2x_3).
\]
The conditional expectations \(P_\Sigma\) and \(P_\Pi\) are the
orthogonal projections onto functions constant on each additive
and multiplicative fiber, respectively. Thus, for a function \(f\),
\[
 (P_\Sigma f)(x)
 =\frac{1}{|\Sigma^{-1}(\Sigma(x))|}
   \sum_{y:\,\Sigma(y)=\Sigma(x)}f(y),
\]
and analogously for \(P_\Pi\). The two projections act on the
same space; they do not represent two independent arithmetic supports.

\begin{lemma}[Two projections and principal angles]
If \(s\in[0,1]\) is a singular value of the pairing between
orthonormal bases of the ranges of two projections \(P,Q\), its nontrivial
block contributes \(s\sqrt{1-s^2}\) to
\(\|[P,Q]\|\).
\end{lemma}
\begin{proof}
On the corresponding principal plane, one can write
\[
 P=\begin{pmatrix}1&0\\0&0\end{pmatrix},\qquad
 Q=\begin{pmatrix}
 s^2&s\sqrt{1-s^2}\\
 s\sqrt{1-s^2}&1-s^2
 \end{pmatrix}.
\]
Their commutator is
\[
 [P,Q]=s\sqrt{1-s^2}
 \begin{pmatrix}0&1\\-1&0\end{pmatrix}.
\]
The final matrix is orthogonal, so its norm is one. Common
or orthogonal blocks have zero commutator. The orthogonal decomposition
into principal planes gives the claim and reduces the total norm to the maximum
of these contributions.
\end{proof}

\begin{proposition}[Exact incompatibility norm]
On \(\ell^2(\mathcal A_3)\),
\begin{equation}
 \|[P_\Sigma,P_\Pi]\|=\frac{\sqrt3}{4}.
 \label{ivloc:norma}
\end{equation}
\end{proposition}
\begin{proof}
Let \(N_{ab}=|\Sigma^{-1}(a)\cap\Pi^{-1}(b)|\). Integer enumeration
of the three symbols yields
\[
 N=9\begin{pmatrix}
 0&1&1&0&1&1&0&1&4\\
 1&0&1&1&0&1&1&0&4\\
 1&0&2&0&1&2&0&0&3\\
 0&1&1&0&1&1&0&1&4\\
 1&0&1&1&0&1&1&0&4\\
 0&0&2&1&0&2&0&1&3\\
 0&1&1&0&1&1&0&1&4\\
 1&0&1&1&0&1&1&0&4\\
 0&1&2&0&0&2&1&0&3
 \end{pmatrix}.
\]
Each row sums to \(81\); the column sums are
\[
 (m_1,\ldots,m_9)=(36,36,108,36,36,108,36,36,297).
\]
The normalized indicator functions of the fibers have pairing
matrix \(C_{ab}=N_{ab}/\sqrt{81m_b}\). Consequently,
\[
 M=CC^{\mathsf T},\qquad
 M_{ac}=\sum_{b=1}^9\frac{N_{ab}N_{cb}}{81m_b}
\]
is a rational matrix whose spectrum consists of the squares of the
required singular values.

Exact multiplication shows that \((1,\ldots,1)\),
\((1,-1,0)^3\), and \((1,1,-2)^3\) are eigenvectors with eigenvalues
\(1\), \(1/4\), and \(7/132\), respectively. The subspace of
vectors supported on \(\{3,6,9\}\) whose sum is zero has dimension
two and eigenvalue \(1/18\). The zero-sum vectors supported on
\(\{1,4,7\}\) or \(\{2,5,8\}\) form a four-dimensional
kernel. These subspaces are independent and their dimensions sum to
nine. The complete spectrum is therefore established:
\[
 \operatorname{spec}(M)=
 \left\{1,\frac14,\frac7{132},\frac1{18},\frac1{18},0,0,0,0\right\}.
\]
The maximum of \(\lambda(1-\lambda)\) on this set is \(3/16\),
attained at \(\lambda=1/4\). The preceding lemma gives
\(\sqrt{3/16}=\sqrt3/4\).
\end{proof}

The mode \((1,-1,0)^3\) is precisely the distribution of values of the
tritic phase selector in the nonadic chart. The scalar norm retains
a measure of incompatibility, but does not by itself preserve the
orientation of that mode or reconstruct the two projectors. That memory
remains in the route state accompanying the scalar reading.

\begin{proposition}[Algebra of the local electronic block]
\label{ivloc:algebra}
On the principal plane attaining \eqref{ivloc:norma}, let
\[
 P=\begin{pmatrix}1&0\\0&0\end{pmatrix},\qquad
 Q=\begin{pmatrix}1/4&\sqrt3/4\\\sqrt3/4&3/4\end{pmatrix},\qquad
 S=2P-I,\qquad J=\frac4{\sqrt3}[P,Q].
\]
Then \(S^2=I\), \(J^2=-I\), \(SJ=-JS\), and the real algebra
generated by \(S,J\) is \(M_2(\mathbb R)\), a realization of
\(\mathrm{Cl}_{1,1}\). Moreover,
\[
 n_\pm=\frac{S\pm J}{2},\qquad
 n_\pm^2=0,\qquad n_+n_-+n_-n_+=I.
\]
\end{proposition}
\begin{proof}
We have \(S=\operatorname{diag}(1,-1)\) and
\(J=\bigl(\begin{smallmatrix}0&1\\-1&0\end{smallmatrix}\bigr)\).
Multiplication verifies the three relations. The matrices
\(I,S,J,SJ\) are linearly independent and form a basis of
\(M_2(\mathbb R)\), proving the claim about the algebra.
Finally, \((S\pm J)^2=S^2+J^2\pm(SJ+JS)=0\), and the sum of the
cross products is \((S^2-J^2)/2=I\).
\end{proof}

In this block, an elliptic direction, a hyperbolic direction, and two
parabolic null directions coexist. The result describes a local
matrix representation of the arithmetic of the two sheets. It does not identify all of APP
with \(M_2(\mathbb R)\), or with a quantum Latin square; such
claims would require other domains and other relations. Nor does the
exponential belong exclusively to the parabolic regime: it can
be applied to any of the generators of this block.


Proposition~\ref{ivloc:algebra} establishes that the pair of projections of the APP sheets
produces a normalized commutator whose square is \(-I\). The same block
contains an involution of square \(I\) and two nilpotent directions.
The joint realization is proved for that concrete pair of projections;
the preceding classification explains why the three quadratic relations
can appear within one matrix algebra without being identified with one another.

In the exceptional prolongation, the Paley matrix reduced modulo \(3\)
satisfies \(A_W^2=-I_6\). This relation endows the ternary space with
a structure of dimension \(3\) over \(\mathbb F_9\). Its connection to
a real operator root is expressed through the integral presentation
\[
 \mathbb Z[u]/(u^2+1).
\]
Base changes to \(\mathbb F_3\) and \(\mathbb R\) produce representations
of the same relation. Each retains its characteristic, dimension, and
incidence information. Developing the respective operators establishes
the connection, not equality of their cardinalities.

Norm-one conservation consequently describes the quadratic realizations
of evolutions. Conservation of quotients, orientation, and trajectory
belongs to the transported arithmetic state. Both properties participate
in the construction, with different functions: one determines the identities
of the representations; the other makes it possible to compose, reconstruct,
and prolong the operations that give rise to them.


\begin{figure}[htbp]
\centering
\begin{tikzpicture}[scale=.88,>=Stealth,font=\small]
\foreach \c/\titulo in {0/{Elliptic},4.6/{Parabolic},9.2/{Hyperbolic}}{
\begin{scope}[xshift=\c cm]
\draw[->,gray] (-1.65,0)--(1.7,0) node[right] {\(u\)};
\draw[->,gray] (0,-1.35)--(0,1.5) node[above] {\(v\)};
\node at (0,2.05) {\titulo};
\end{scope}}
\draw[thick] (0,0) circle[radius=1];
\draw[->,thick] (1,0) arc[start angle=0,end angle=60,radius=1];
\node at (0,-1.85) {\(u^2+v^2=1\)};
\node at (0,-2.35) {\(J^2=-I\)};
\draw[thick] (3.6,-1.2)--(3.6,1.2);
\draw[->,thick] (5.6,-1.2)--(5.6,1.2);
\node at (4.6,-1.85) {\(u^2=1\)};
\node at (4.6,-2.35) {\(J^2=0\)};
\draw[thick,domain=-.9:.9,samples=50] plot ({9.2+cosh(\x)},{sinh(\x)});
\draw[thick,domain=-.9:.9,samples=50] plot ({9.2-cosh(\x)},{sinh(\x)});
\node at (9.2,-1.85) {\(u^2-v^2=1\)};
\node at (9.2,-2.35) {\(J^2=I\)};
\end{tikzpicture}
\caption{Quadratic realizations of TRIT. For
\(J_\tau=\left(\begin{smallmatrix}0&-\tau\\1&0\end{smallmatrix}\right)\),
the exponential preserves \(u^2+\tau v^2\). Parabolic degeneration
preserves \(u\) and transports \(v\) linearly. The three diagrams represent
orbits of quadratic forms; their geometric type corresponds to the
algebraic relation of the generator.}
\label{fig:trit-regimenes}
\end{figure}


\newpage
\subsection{TPK: Topological Phase Kernel}

The dynamics uses the phase selector
\[
 M_{\rm ph}(r)=
 \begin{cases}
 +1,&r\in\{1,4,7\},\\
 -1,&r\in\{2,5,8\},\\
 0,&r\in\{3,6,9\}.
 \end{cases}
\]
Its value determines which evaluation operates. Cardinal transport moves the corresponding cursor; the record preserves the previous value, the new quotient, the orientation, and the transition. Selection and transport are therefore different functions: a ternary sign is not a displacement on the torus.

We shall denote by \(\widehat X\) the space of arithmetic states with trajectory memory. An element contains the positions and directions of both cursors, the active mode, the TRIT signature, the phase, the residues and quotients of both evaluations, the carries, the constructed prefix, and the boundary data. When a refinement cylinder is introduced, it forms part of the state. This notation avoids replacing the object by a reduced version that retains only the visible symbol.

\begin{definition}[TPK update]
Let \(\mathsf{Sel}_t\) be mode selection through \(M_{\rm ph}\), \(\mathsf{Tra}_t\) the lifted cardinal transport, and \(\mathsf{Mem}_t\) the record update. The transition is
\begin{equation}
 \mathcal U_t=\mathsf{Mem}_t\circ\mathsf{Tra}_t\circ\mathsf{Sel}_t.
 \label{eq:actualizacion}
\end{equation}
Each application preserves the fields it does not modify and updates the remaining fields by Euclidean division and trajectory concatenation.
\end{definition}

The notation \eqref{eq:actualizacion} expresses the order of composition; it does not replace the effective rules. The following section specifies its finite two-cursor realization completely. Prolongation will then be obtained not by erasing its memory and repeating the first emission, but by transporting the prefixes and their compatible boundaries. Thus APP supplies the evaluations, TRIT distinguishes their regimes, and TPK composes their evolution.

% Expository recovery from the TPK kernel's source owners.
% Provenance and preservation controls are documented in technical/TPK_INTEGRACION.md.
% This file expands the TPK section; it does not replace extension.tex or generacion.tex.

\subsubsection{Dynamic state, observables, and operational dependencies}
\label{tpk-int:estado}

The Topological Phase Kernel organizes the evolution of the two APP sheets
with the orientation determined by TRIT. Its definition comprises the state
on which it acts, the transition rules, and the relations that make
prolongations compatible. Equation \eqref{eq:actualizacion} acquires operational
content when the coordinates received and modified by each factor are specified.
Position belongs to the toroidal support; phase belongs to the calendar;
orientation determines transport; quotients record the arithmetic evaluation
before reduction. These coordinates are related through the transition and
retain different mathematical functions.

The observable projection of a state is
\begin{equation}
 q=(x_+,d_+,x_\times,d_\times,\theta)
 \in\mathcal Q_{\rm obs}
 :=(X_9\times\mathcal D)^2\times\Theta_{108},
 \qquad \mathcal D=\{N,E,S,O\},\quad
 \Theta_{108}=\ZZ/108\ZZ.
 \label{tpk-int:qobs}
\end{equation}
The subscripts distinguish the additive and multiplicative cursors. The
fibers of complete states over \(q\) preserve the ordered trajectory, the
integer evaluations of both sheets, their residues and quotients, and the
orientation and boundary records. A precision prolongation adds the prefix
and the compatible cylinder. Thus two states with the same projection
\(q\) may retain different records.

Operational dependencies have three levels. At the local level, the selector
activates a sheet, transport changes its position, and the update computes
the new arithmetic coordinates. At the trajectory level, these transitions
compose and retain the records of preceding steps. At the refinement level,
extension relations determine which continuation preserves the predecessor's
prefix, signature, and boundary. The finite calendar describes the phase
of these operations; holonomy describes their composition on the history.

\subsubsection{Tritic selection and lifted cardinal transport}
\label{tpk-int:seleccion}

Let \(\bar\theta\in\{0,\ldots,107\}\) be the calendar representative.
The operating phase and the dodecaphase sector are
\begin{equation}
 r(\theta)=1+(\bar\theta\bmod9),\qquad
 \beta(\theta)=\left[\left\lfloor\frac{\bar\theta}{9}\right\rfloor\right]_{12}.
 \label{tpk-int:fase}
\end{equation}
The previously defined selector may be written \(\varepsilon_9=M_{\rm ph}\).
Its value vector in the ordered chart of \(D_9\) is
\[
 m_{\rm ph}=(1,-1,0,1,-1,0,1,-1,0).
\]
The function \(M_{\rm ph}:D_9\to\TRIT\) and this vector are two
presentations of the same selector, with their respective domains.

Define the activation indicators
\[
 a_+(\theta)=\mathbf1_{\{1,4,7\}}(r(\theta)),\qquad
 a_\times(\theta)=\mathbf1_{\{2,5,8\}}(r(\theta)).
\]
Cursor transport is then
\begin{equation}
 x_+'=x_++a_+(\theta)\nu_{d_+},\qquad
 x_\times'=x_\times+a_\times(\theta)\nu_{d_\times}
 \quad\text{in }X_9.
 \label{tpk-int:cursores}
\end{equation}
In phases \(3,6,9\), both positions remain unchanged; the neutral event
forms part of the chronology and its record. Cardinal directions determine
the vectors \(\nu_d\), whereas the phase signature determines which
vector is applied.

The integer lift preserves boundary crossing. If
\(s:X_9\to D_9^2\) chooses positive representatives and
\(x'=x+a\nu_d\), with \(a\in\{0,1\}\), define
\begin{equation}
 b_d(x,a)=\frac{s(x)+a\nu_d-s(x')}{9}\in\ZZ^2.
 \label{tpk-int:borde}
\end{equation}
Divisibility by \(9\) follows from equality in \(X_9\).
If \(z\in\ZZ^2\) records previous crossings, the update is
\(z'=z+b_d(x,a)\). Thus
\[
 s(x')+9z'=s(x)+9z+a\nu_d.
\]
The integer transport data remain reconstructible after the position
has been reduced on the torus.

\begin{proposition}[Conservation of lifted displacement]
\label{tpk-int:prop-desplazamiento}
For a cardinal trajectory of \(n\) steps, with activation indicators
\(a_t\),
\[
 s(x_n)+9z_n=s(x_0)+9z_0+
 \sum_{t=0}^{n-1}a_t\nu_{d_t}.
\]
The same equality is preserved under concatenation of two compatible trajectories.
\end{proposition}
\begin{proof}
The one-step identity is \eqref{tpk-int:borde}. Summing it over the
\(n\) steps cancels the intermediate coordinates. Splitting the summation
interval into two segments produces exactly the concatenation law.
For a closed trajectory, the difference \(z_n-z_0\) is the previously
defined winding number.
\end{proof}

The same principle governs APP evaluations. If a sheet takes the integer
value \(N_t\), its record simultaneously preserves
\(N_t=R_t+9q_t\). The difference between two steps satisfies
\[
 N_{t+1}-N_t=(R_{t+1}-R_t)+9(q_{t+1}-q_t).
\]
Spatial-crossing memory \(z_t\) and arithmetic quotient \(q_t\)
have different definitions; preserving them jointly maintains the
geometric and arithmetic provenance of every transition.

\subsubsection{Two-cursor chronology and composite returns}
\label{tpk-int:cronologia}

The cardinal involution \(\jmath\) interchanges \(N\) with \(S\) and
\(E\) with \(O\). After applying \eqref{tpk-int:cursores}, both
directions are reversed at the end of a block of \(27\) steps.
Writing
\[
 h(\theta)=\mathbf1_{\{0,27,54,81\}}
             ((\bar\theta+1)\bmod108),
\]
the complete rule on the observable projection is
\begin{equation}
 F_{\rm obs}(q)=
 \bigl(x_+',\jmath^{h(\theta)}d_+,
       x_\times',\jmath^{h(\theta)}d_\times,
       \theta+1\bigr).
 \label{tpk-int:fobs}
\end{equation}
The reading that forms the emission uses the positions before transport,
according to the explicit recurrence in Section~\ref{sec:extension}.
The observable update and emission therefore compose in a fixed order.

\begin{proposition}[Reversibility and period of the observable calendar]
\label{tpk-int:periodo}
The map \(F_{\rm obs}\) is a permutation of
\(\mathcal Q_{\rm obs}\) and satisfies \(F_{\rm obs}^{108}=I\).
Within a block of \(54\) steps, positions and directions are restored;
the coordinate in \(\Theta_{108}\) advances by \(54\).
\end{proposition}
\begin{proof}
Given the final state, the preceding phase is first recovered by subtracting
one in \(\Theta_{108}\). That phase determines whether directions were
reversed and which cursor moved. Applying the corresponding involution
and subtracting the cardinal vector uniquely recovers the initial state.

Each interval of \(27\) steps contains nine activations of each cursor.
The following block preserves the activations and reverses the directions,
so the integer displacements cancel over \(54\) steps. The rule has
period \(54\) in its position and orientation increments; the same
cancellation holds for any calendar origin. After \(108\) steps,
\(\theta\) is also restored.
\end{proof}

We reserve
\[
 \mathcal K_{27}=F_{\rm obs}^{27}:
 \mathcal Q_{\rm obs}\longrightarrow\mathcal Q_{\rm obs}
\]
for the iterate of \(27\) transitions. Its fourth iterate is the
identity on the observable state defined in \eqref{tpk-int:qobs}.
Projecting history onto the calendar makes this finite check possible.
The trajectory record, its evaluations, and the new prefixes remain in
the dynamic state on which \(\mathcal U_t\) acts; observable equality
does not reinitialize them.

\subsubsection{Affine memory transport and composition of states}
\label{tpk-int:memoria}

The observable projection admits fibers of arithmetic and geometric data.
Over a configuration \(q\), we write a state in the form
\[
 x=(q,o,h,c,\sigma,C,b,\lambda).
\]
Here \(o\) preserves orientation and sheet; \(h\), residues and
quotients; \(c\), additive transport memory; \(\sigma\), the boundary
signature; \(C\), the precision cylinder; \(b\), its boundary label;
and \(\lambda\), the recorded trajectory. The signature is determined
by a map
\[
 \operatorname{Sig}_q:
 \mathsf H_q\times\mathsf C_q\times\mathsf{Hist}_q
 \longrightarrow\mathsf S_q,
 \qquad \sigma=\operatorname{Sig}_q(h,c,\lambda).
\]
The symbols \(\mathsf H_q,\mathsf C_q,\mathsf{Hist}_q\), and
\(\mathsf S_q\) denote, respectively, the sets of residue--quotient
data, the abelian transport-memory group, the recorded trajectories,
and the boundary signatures. Realizations of this signature are specified
component by component in the regional and terminal constructions.
This dependence prevents treating the signature as a free parameter
once the data producing it are fixed.

The elementary residue--quotient chart is explicit. For
\(n\in\ZZ\), let
\[
 r=1+((n-1)\bmod9),\qquad u=\frac{n-r}{9}.
\]
Then \(n=r+9u\), with \(r\in D_9\) and \(u\in\ZZ\).
Uniqueness follows because two representatives in \(D_9\) congruent
modulo \(9\) are equal. This chart supplies local integer reconstruction.
The precision towers additionally incorporate the truncation maps and
cylinders developed in \ref{sec:extension} and \ref{sec:generacion}.

Let \(\gamma:q\to q'\) be an observable trajectory. Fiber transport
is expressed through maps
\[
 \mathsf H(\gamma):\mathsf H_q\to\mathsf H_{q'},\qquad
 \mathsf C(\gamma):\mathsf C_q\to\mathsf C_{q'},\qquad
 \mathsf{Hist}(\gamma):\mathsf{Hist}_q\to\mathsf{Hist}_{q'}.
\]
The maps respect identities and composition, and
\(\mathsf C(\gamma)\) is an abelian-group homomorphism.
The increment \(k_{\rm tr}(\gamma)\in\mathsf C_{q'}\) satisfies
\begin{equation}
 \begin{split}
 k_{\rm tr}(\operatorname{id}_q)&=0,\\
 k_{\rm tr}(\gamma_2\circ\gamma_1)
 &=k_{\rm tr}(\gamma_2)+
   \mathsf C(\gamma_2)k_{\rm tr}(\gamma_1).
 \end{split}
 \label{tpk-int:cociclo}
\end{equation}
Here \(k_{\rm tr}\) denotes the transport cocycle;
the rational coordinate \(\kappa\) of record \(K\) retains
its notation in the dodecaphase-closure chapter.

Transportable coordinates are updated by
\begin{equation}
 \begin{aligned}
 h'&=\mathsf H(\gamma)h,&
 c'&=\mathsf C(\gamma)c+k_{\rm tr}(\gamma),\\
 \lambda'&=\gamma\circ\lambda,&
 \sigma'&=\operatorname{Sig}_{q'}(h',c',\lambda').
 \end{aligned}
 \label{tpk-int:accion-memoria}
\end{equation}
The second equation is an affine action: the linear term transports
previous memory and the cocycle adds the increment produced by the
trajectory. In the spatial-crossing realization for each cursor separately,
\(\mathsf C(\gamma)=I\) and \(k_{\rm tr}(\gamma)\) is the sum
of the vectors in \eqref{tpk-int:borde}. This concrete instance connects
the abstract law with the cardinal operations already defined.

\begin{proposition}[Compatibility of the update with concatenation]
\label{tpk-int:composicion-memoria}
Updating by \eqref{tpk-int:accion-memoria} along
\(\gamma_1\), followed by updating along \(\gamma_2\),
coincides with updating along \(\gamma_2\circ\gamma_1\).
\end{proposition}
\begin{proof}
For the affine coordinate, two updates give
\[
 c''=\mathsf C(\gamma_2)\mathsf C(\gamma_1)c
       +\mathsf C(\gamma_2)k_{\rm tr}(\gamma_1)
       +k_{\rm tr}(\gamma_2).
\]
Functoriality of \(\mathsf C\) and
\eqref{tpk-int:cociclo} turn this expression into
\(\mathsf C(\gamma_2\circ\gamma_1)c+
k_{\rm tr}(\gamma_2\circ\gamma_1)\).
The statement for \(h\) follows from composition of
\(\mathsf H\), and that for \(\lambda\) from associativity
of trajectories. The final signature agrees because it is evaluated
on the same three reconstructed coordinates.
\end{proof}

The cylinder and its boundary complete this law through the relations
\(\operatorname{Ext}_\gamma\), defined in
\ref{tpk-int:bidireccionalidad}. Their elements determine which of
the preceding updates preserve an admissible prolongation.
States, together with the triples \((\gamma,x,x')\) admitted by
these relations, form a category over observable trajectories:
\[
 (\gamma_2,x',x'')\circ(\gamma_1,x,x')
     =(\gamma_2\circ\gamma_1,x,x'').
\]
The identity of \(x\) is \((\operatorname{id}_q,x,x)\).
Composition preserves arithmetic transport and boundary prolongability
simultaneously. This is the structure in which returns with memory
and successive refinements take place.

\subsubsection{Dodecaphase record, phase, and accumulated memory}
\label{tpk-int:dodecafase}

The \(108\) transitions are distributed over the ordered windows
\[
 E_m=\{9m-8,\ldots,9m\},\qquad 1\leq m\leq12.
\]
The record preserves the phase origin, orientation, and data
produced during each window. Its map is
\[
 \mathcal R_{12}(x)=
 \bigl((E_m,\tau_m,\sigma_m,\kappa_m,\Lambda_m)\bigr)_{m=1}^{12},
\]
where \(\tau_m\) is the window's subroute, \(\sigma_m\) its
orientation, \(\kappa_m\) the integer boundary increment, and
\(\Lambda_m\) the local incidence record. We adopt the chart
developed in \ref{subsec:k-registro-integral}; its window indices
correspond to \(m=\beta+1\).
The domain comprises compatible TPK histories, and the codomain
their oriented temporal records. Composition of segments and their
increments is governed by \eqref{tpk-int:accion-memoria}.
By contrast, the group \(C_{12}\) describes translation of the
window index; \(\Theta_{108}\) preserves the step calendar.
This distinction specifies what is transformed and what is observed
when a cycle is completed.

In coordinates \(\theta=9b+r\), with \(0\leq r<9\), advancing
\(9\) steps preserves \(r\) and increases \(b\) by one.
Addition of phases introduces the quotient
\[
 c_0(r,r')=\left\lfloor\frac{r+r'}9\right\rfloor,
\]
whose cocycle identity ensures associativity of composition.
The finite calendar preserves \([b]_{12}\); a record of turns
preserves \(b\in\ZZ\), or \(b\in\mathbb N_0\) for a
history with fixed origin. Section~\ref{sec:extension} proves
the nonsplit exact sequence \eqref{eq:extension108}. Its function
is to describe the relation between these coordinates before they
are used to construct the signed state.

The signed record receives phase and memory data through the
composition developed in Section~\ref{sec:k-registro}:
\begin{equation}
 R_{\rm sgn}=\Pi_H\circ\mathcal E_{108}^{90,120}
                         \circ\mathcal R_{12}.
 \label{tpk-int:registro-causal}
\end{equation}
Its domain is the terminal state with memory; its output belongs
to \(\ZZ^{12}\). The reversible transformation between this signed
state and \(K\), together with the rational reading of \(K\),
is proved in the corresponding chapter. Record operations precede
these transformations. Reconstructing a record from its cyclic
differences verifies its recoverability; its dynamic provenance
must be preserved through the events that produced it.

For the encoded trace \((d_t)\) retained by the record,
the event reader counts codes \(\{2,4,6,8\}\) in each window
with orientation \((+,+,+,-,-,-,+,+,+,-,-,-)\).
These codes and the cardinal symbols of \(\mathcal D\)
are different record coordinates. Chapter~\ref{sec:k-registro}
develops this reader, the integral decomposition \(1+5+6\),
the congruences characterizing its image, and its exact inverse.
The inversion proof establishes which data this chart preserves;
the definition of \(\mathcal R_{12}\) specifies the oriented
history on which it acts.

\subsubsection{Phase incidence and internal channel construction}
\label{tpk-int:incidencia}

The selector determines three subsets of \(D_9\):
\[
 H_+=\{1,4,7\},\qquad H_-=\{2,5,8\},\qquad H_0=\{3,6,9\}.
\]
Let \(H_{\rm act}=H_+\sqcup H_-\) be the set of active phases
and \(O_{\rm ph}=D_9\setminus\{9\}\) the chart preceding
the marked return. Denote by \(P_2(H_{\rm act})\) the set of
unordered pairs of active phases. Internal incidence is given by
\begin{equation}
 \mathcal F^{\rm ph}_6=H_{\rm act}\times P_2(H_{\rm act}),\qquad
 \mathcal F^{\rm ph}_8=O_{\rm ph}\times P_2(H_{\rm act}),\qquad
 \Theta_{\rm ph}=D_9\times P_2(H_{\rm act}).
 \label{tpk-int:fases90120}
\end{equation}
The indices describe the number of positions in the first component.
The construction uses the phases and calendar origin already fixed;
subsequent exceptional realizations will transport these incidences
to their own domains.

\begin{proposition}[Cardinalities and oriented phase incidence]
\label{tpk-int:cardinales}
The sets in \eqref{tpk-int:fases90120} have cardinalities
\(90,120,135\), respectively. If
\[
 \partial_{\rm ph}=\Theta_{\rm ph}\times\{-1,+1\},\qquad
 E_{\rm ph}=\partial_{\rm ph}\sqcup\{\infty\},\qquad
 L_{\rm ph}=H_0\times P_2(H_{\rm act}),
\]
then
\[
 |\partial_{\rm ph}|=270,\quad |E_{\rm ph}|=271,\quad
 |L_{\rm ph}|=45,\quad
 |\partial_{\rm ph}\sqcup L_{\rm ph}|=315.
\]
Simultaneously translating the phase origin and marked return preserves
all these cardinalities and inclusions.
\end{proposition}
\begin{proof}
There are six active phases, so
\(|P_2(H_{\rm act})|=\binom62=15\). Multiplying by
\(6,8,9\) gives \(90,120,135\). Orientation doubles \(135\);
the return mark adds one element; the three neutral phases give
\(3\cdot15=45\). A calendar translation transports each factor
by a bijection and carries the marked point to its image. Products
and disjoint unions thus preserve the incidence.
\end{proof}

Here the cardinality \(271\) has a phase-incidence realization.
The cubic completion in Section~\ref{sec:extension} has its own
construction of the corona \(1000-729\). The correspondence
between the two realizations must preserve the maps and marks
identifying them, in addition to their cardinality. Likewise,
the selector \(\varepsilon_9\), the incidence \((90,120)\),
and the extractor \(\mathcal E_{108}^{90,120}\) retain the
domains and operations just distinguished.

\subsubsection{Trajectory reversal and conjugate prolongations}
\label{tpk-int:bidireccionalidad}

Bidirectionality begins with an operation on trajectories before scalar
values are assigned to them. For
\(\gamma=(x_0;d_1,\ldots,d_n)\), with endpoint \(x_n\), define
\[
 \operatorname{rev}\gamma
   =(x_n;\jmath d_n,\ldots,\jmath d_1).
\]
The reversed trajectory starts at the endpoint of the original.
Composition satisfies
\[
 \operatorname{rev}^2\gamma=\gamma,\qquad
 \operatorname{rev}(\gamma_2\circ\gamma_1)
   =\operatorname{rev}\gamma_1\circ\operatorname{rev}\gamma_2,
\]
and integer displacement changes sign. These identities follow by
reversing the order of steps and replacing each cardinal vector
with its opposite. For a loop, it follows that
\(\operatorname{wind}(\operatorname{rev}\gamma)
=-\operatorname{wind}(\gamma)\).

On records, reversal also requires transporting the order of components
and the orientation of events. Route reversal and signature conjugation
are distinguishable operations: their composition is specified when
applied to a concrete record. In particular, the regional involution
interchanging closure and propagation components preserves the self-scaling
axis; its formula is given in Section~\ref{mono-ampl:regiones}.
This regional operation acts on a different domain from cardinal
trajectory reversal.

Extensions over a route \(r:q\to q'\) are expressed by
relations between fibers of compatible states. If \(\mathcal F_q\)
denotes the fiber of complete states over \(q\), we require
\[
 \operatorname{Ext}_r\subseteq\mathcal F_q\times\mathcal F_{q'},
 \qquad \Delta_{\mathcal F_q}\subseteq
        \operatorname{Ext}_{\operatorname{id}_q},
\]
\[
 \operatorname{Ext}_{r_2}\circ\operatorname{Ext}_{r_1}
       \subseteq\operatorname{Ext}_{r_2\circ r_1}.
\]
Their composition preserves the evaluations of both sheets, orientation,
quotients, prefix, and boundary. Route reversal does not by itself imply
that every extension relation is a bijection. Conjugate prolongation
is defined on the domain where transport of these data satisfies
the corresponding compatibilities.

\subsubsection{Nonadic holonomy and continuity of refinement}
\label{tpk-int:holonomia}

A phase turn composes transitions on the state with memory:
\begin{equation}
 \Gamma_{9,k}=\mathcal U_{k+8}\circ\cdots\circ\mathcal U_k.
 \label{tpk-int:gamma}
\end{equation}
On an individualized history, this functional composition realizes
the holonomy relation; on the full domain we retain
\(\Gamma_{9,k}\subseteq\widehat X_k\times\widehat X_{k+9}\),
with the admissible prolongations developed later.
The cyclic base has nine phases, while the fiber preserves trajectory,
quotients, and boundary. Monodromy is the action of one turn of this
base on the fibers; holonomy \(\Gamma_{9,k}\) expresses its
realization through the effective TPK transitions.

If \(g\) observes phase and \(M\) the number of retained turns,
\[
 g(\Gamma_{9,k}x)=g(x),\qquad M(\Gamma_{9,k}x)=M(x)+1.
\]
The first identity expresses phase closure and the second identifies
the new temporal depth. Coinductive prolongation also preserves its
cylinders and prefixes. Restrictions between depths compose an inverse
system; a compatible section gathers its antecedents without replacing
them by the latest observation. This arithmetic of histories and its
unity are developed in \ref{hist:seccion}--\ref{hist:doble-varianza};
the regions and their readers are constructed in Section~\ref{sec:generacion}.

Phase transfer on emissions already constructed has a different domain.
Let \(\mathcal A_+\) and \(\mathcal A_\times\) be the
additive and multiplicative signature alphabets, with cardinalities
\(18\) and \(26\). Their product identifies the catalog \(\mathcal U_6\).
For \(f:\mathcal A_+\times\mathcal A_\times\to\RR\),
the sheet averages are
\[
 (E_+f)(s,e)=\frac1{18}\sum_{s'\in\mathcal A_+}f(s',e),\qquad
 (E_\times f)(s,e)=\frac1{26}\sum_{e'\in\mathcal A_\times}f(s,e').
\]
These emissions and their cardinalities are constructed through the
recurrence in Section~\ref{sec:extension}. Once they have been obtained,
the transfer operator is defined on
\(\mathcal U_6\times C_9\) by
\[
 P_+=E_+,\quad P_-=E_\times,\quad P_0=I,\qquad
 \mathcal K_{\rm ph}((u,r),(v,r+1))=P_{M_{\rm ph}(r)}(u,v),
\]
with zero weight for all other phase transitions. The index
\(r\) uses the positive representative of the cyclic class.
The phase order \((+1,-1,0)^3\) then produces the renewal
\[
 \mathcal K_{\rm ph}^{9}\big|_r=E_+E_\times.
\]
Indeed, both averages are idempotent, commute, and average different
factors; their composition assigns weight \(1/(18\cdot26)=1/468\)
to each emission. This equality characterizes a subsequent transfer
observable. History transport remains given by \eqref{tpk-int:gamma},
with its memory and boundary.

The resulting hierarchy establishes the function of the following chapters.
Finite emission constructs the domain and regions; extension relations
preserve their prefixes; the nonadic connection prolongs them; the
dodecaphase record determines the data entering \(K\) and the
closure of \(\alpha\). Numerical and incidence realizations receive
this preceding structure. This order allows every consequence to be
presented in its own chapter while preserving, from the formal kernel
onwards, the objects that make it possible.


\subsection{Accumulated memory and transport resolvents}
\label{subsec:memoria-resolvente}

Calendar return does not remove increments produced during the journey.
To express this distinction in operator form, an accumulated trajectory-event
record is constructed first, followed by its generating series. Reduction
of memory variables will be performed on this transport, also preserving
the source and output reader.

\subsubsection{Oriented events and record update}

Let \((d_t)\) be the sequence of integer direction codes retained by
a TPK trajectory. The code \(d_t\) and the cursor's cardinal direction
are different coordinates; no identification of their alphabets is assumed.
We number the windows starting from zero:
\[
 I_m=\{9m+1,\ldots,9m+9\},\qquad m\geq0.
\]
In the first cycle, \(I_m\) is window \(E_{m+1}\) of the
dodecaphase record. Its orientation is
\[
\sigma=(1,1,1,-1,-1,-1,1,1,1,-1,-1,-1).
\]
For a prolonged history, \(\sigma_m\in\{-1,1\}\) is the sign
retained in the corresponding window. The signed event is
\begin{equation}
 a_m=\sigma_m\sum_{t\in I_m}
       \mathbf1_{\{2,4,6,8\}}(d_t),\qquad |a_m|\leq9.
 \label{mem:evento}
\end{equation}
Each summand is incorporated when its transition occurs. This definition
counts trace events; it does not select directions through a value
that is to be obtained later.

Let \(\mathbf e_0,\ldots,\mathbf e_{11}\) be the basis of \(\ZZ^{12}\)
and \(S\mathbf e_j=\mathbf e_{j+1\bmod12}\). Write
\(R=S^{-1}\). The accumulated record in fixed coordinates satisfies
\[
z_{m+1}=z_m+a_m\mathbf e_{m\bmod12}.
\]
A prefix with no accumulated history starts at \(z_0=0\); at a different
origin, \(z_0\) retains the preceding record.
The frame accompanying the calendar is \(y_m=S^{-m}z_m\).

\begin{proposition}[Update and return with memory]
\label{mem:prop-actualizacion}
The record in the moving frame satisfies
\begin{equation}
 y_{m+1}=Ry_m+\eta_m,\qquad
 \eta_m=a_mR\mathbf e_0.
 \label{mem:actualizacion}
\end{equation}
For every \(n\geq1\),
\begin{equation}
 y_n=R^ny_0+\sum_{j=0}^{n-1}R^{n-1-j}\eta_j.
 \label{mem:iteracion}
\end{equation}
In particular,
\begin{equation}
 y_{12}=y_0+\sum_{j=0}^{11}a_j\mathbf e_j.
 \label{mem:retorno12}
\end{equation}
\end{proposition}
\begin{proof}
Multiplying the update of \(z_m\) by \(S^{-(m+1)}\) gives
\(y_{m+1}=S^{-1}y_m+a_mS^{-1}\mathbf e_0\), since
\(S^{-m}\mathbf e_{m\bmod12}=\mathbf e_0\).
Formula \eqref{mem:iteracion} follows by induction.
For \(n=12\), \(R^{12}=I\) and
\(R^{11-j}\eta_j=a_jR^{12-j}\mathbf e_0=a_j\mathbf e_j\),
which proves \eqref{mem:retorno12}.
\end{proof}

Thus twelve windows restore the frame and preserve the event vector.
The count alone does not recover the order of all steps within each
window; that order remains in the recorded trajectory.

To specify the balance, define
\[
 D_3=I-S^3,\quad D_4=I-S^4,\quad
 B=D_3^{\mathsf T}D_3+D_4^{\mathsf T}D_4,
\]
\[
 \mathcal M(y)=\tfrac12y^{\mathsf T}By,
 \qquad q(y)=\sum_{j=0}^{11}y_j.
\]
The first is a quadratic form of the record; the second records the sum
of its coordinates. They are not identified here with physical energy or charge.

\begin{proposition}[Event-induced balance]
\label{mem:prop-balance}
Update \eqref{mem:actualizacion} satisfies
\begin{equation}
 \begin{aligned}
 \mathcal M(y_{m+1})-\mathcal M(y_m)
   &=a_m(By_m)_0+2a_m^2,\\
 q(y_{m+1})-q(y_m)&=a_m.
 \end{aligned}
 \label{mem:balance}
\end{equation}
\end{proposition}
\begin{proof}
Expanding the products gives
\(B=4I-S^3-S^{-3}-S^4-S^{-4}\); hence
\(R^{\mathsf T}BR=B\) and \(B_{00}=4\).
Since \(y_{m+1}=R(y_m+a_m\mathbf e_0)\), the quadratic difference
is \(a_m\mathbf e_0^{\mathsf T}By_m+\tfrac12a_m^2B_{00}\).
The coordinate sum is invariant under permutation \(R\),
giving the second equality.
\end{proof}

\Needspace{20\baselineskip}
\subsubsection{Generating series and uniform memory control}

The series preserves all increments, not only their sum at the end of
a cycle. With a formal variable \(u\), let
\[
 Y(u)=\sum_{m\geq0}y_mu^m,\qquad
 H_\eta(u)=\sum_{m\geq0}\eta_mu^m.
\]

\begin{proposition}[Resolvent and tail bound]
\label{mem:prop-serie}
The formal identity
\begin{equation}
 Y(u)=(I-uR)^{-1}\bigl(y_0+uH_\eta(u)\bigr)
 \label{mem:serie}
\end{equation}
holds. Both series converge for \(|u|<1\). If \(\ell:\RR^{12}\to\RR\)
is linear, \(L=\|\ell\|_{1\to\RR}\), \(M_0=\|y_0\|_1\), and
\(0<r<1\), then, for every integer \(N\geq0\) and \(|u|\leq r\),
\begin{equation}
 \left|\sum_{m>N}\ell(y_m)u^m\right|
 \leq Lr^{N+1}\left(
 \frac{M_0+9(N+1)}{1-r}+\frac{9r}{(1-r)^2}\right).
 \label{mem:cota}
\end{equation}
The derivative series converges uniformly on every disk
\(|u|\leq r<1\).
\end{proposition}
\begin{proof}
Multiplying \eqref{mem:actualizacion} by \(u^{m+1}\) and summing gives
\(Y-y_0=uRY+uH_\eta\). The constant term of \(I-uR\) is
invertible, so its formal inverse is \(\sum_{k\geq0}u^kR^k\).
Moreover, \(R\) preserves the \(\ell^1\) norm and
\(\|\eta_m\|_1\leq9\), whence
\(\|y_m\|_1\leq M_0+9m\).
For the tail, sum the majorants
\(L(M_0+9m)r^m\), using
\[
 \sum_{m>N}r^m=\frac{r^{N+1}}{1-r},\qquad
 \sum_{m>N}mr^m=r^{N+1}
 \left(\frac{N+1}{1-r}+\frac r{(1-r)^2}\right).
\]
Finally, \(\sum_{m\geq1}Lm(M_0+9m)r^{m-1}<\infty\)
majorizes the derivative series and proves its uniform convergence.
\end{proof}

The variable \(u\) counts windows of nine transitions. Identifying it
with another reader's variable requires preserving that clock and the
map relating the two readings.

\subsubsection{Memory elimination: operator, source, and reader}

On a domain where the update is fixed as \(U:X\to X\),
the free vector space \(V=\QQ^{(X)}\) on the states admits the operator
\(T\delta_x=\delta_{U(x)}\). The clock is included in \(x\) if
the transition depends on it. This linear extension preserves distinct
states even when they share an observable phase.

Consider a decomposition \(V=V_0\oplus V_1\), where \(V_1\)
is the auxiliary sector to be eliminated, and write
\[
 T=\begin{pmatrix}A&B_1\\C&D\end{pmatrix},\qquad
 v=\binom{v_0}{v_1},\qquad \ell=(\ell_0,\ell_1).
\]
Here \(A:V_0\to V_0\), \(B_1:V_1\to V_0\),
\(C:V_0\to V_1\), and \(D:V_1\to V_1\).
The following identities are formal; they also hold wherever the
resolvents converge.

\begin{proposition}[Schur reduction with source and reader]
\label{mem:prop-schur}
Let
\[
 D_u=(I-uD)^{-1},\qquad
 \mathscr F(u)=I-uA-u^2B_1D_uC.
\]
The retained component of \((I-uT)^{-1}v\) is
\begin{equation}
 w_0=\mathscr F(u)^{-1}\bigl(v_0+uB_1D_uv_1\bigr),
 \qquad w_1=D_u(v_1+uCw_0).
 \label{mem:schur-fuente}
\end{equation}
Its complete reading satisfies
\begin{equation}
 \begin{split}
 \ell(I-uT)^{-1}v
 ={}&(\ell_0+u\ell_1D_uC)
       \mathscr F(u)^{-1}(v_0+uB_1D_uv_1)\\
    &+\ell_1D_uv_1.
 \end{split}
 \label{mem:schur-lector}
\end{equation}
\end{proposition}
\begin{proof}
The two equations of \((I-uT)w=v\) are
\(w_0=v_0+uAw_0+uB_1w_1\) and
\(w_1=v_1+uCw_0+uDw_1\).
Solving the second and substituting it into the first gives
\eqref{mem:schur-fuente}; applying \(\ell_0\) and \(\ell_1\)
to the two components gives \eqref{mem:schur-lector}.
All formal inverses exist because their constant terms are identities.
\end{proof}

The term \(u^2B_1D_uC=\sum_{k\geq0}u^{k+2}B_1D^kC\)
describes departure into memory, \(k\) steps within it, and a return.
By contrast, \(u\ell_1D_uC\) reads memory before that return.
Thus reducing only the operator does not necessarily preserve the
observation: the source and reader must also be transformed.

The distinction is visible in the cycle already constructed. If \(P\)
retains \(\mathbf e_0\), then \(PRP=0\), but
\begin{equation}
 \left\langle\mathbf e_0,(I-uR)^{-1}\mathbf e_0\right\rangle
 =\frac1{1-u^{12}},\qquad
 \left\langle\mathbf e_{11},(I-uR)^{-1}\mathbf e_0\right\rangle
 =\frac u{1-u^{12}}.
 \label{mem:ejemplo-ciclo}
\end{equation}
Indeed, \(R^n\mathbf e_0=\mathbf e_{-n\bmod12}\): the two
series collect, respectively, \(n\equiv0\) and
\(n\equiv1\pmod{12}\). The resolvent of the compression is \(I\),
whereas compression of the resolvent preserves the omitted returns.
This example concerns the twelve-window record, not the entire TPK state.

\subsubsection{Segment composition and coherence of reduction}

Let three composable affine transports be
\(F_i(y)=R_i y+b_i\), with compatible domains and codomains.
The memory produced by the complete journey is
\begin{equation}
 F_3F_2F_1(y)
 =R_3R_2R_1y+b_3+R_3b_2+R_3R_2b_1.
 \label{mem:cociclo-tres}
\end{equation}
The proof is successive substitution of the three affine formulas.
Grouping the first two segments first, or the last two, produces the
same expression; the factors transport each increment from its point
of origin. In \eqref{mem:actualizacion},
\(R_i=R\) and \(b_i\) is the event of its window.

Eliminating several layers preserves a similar coherence.
To make it explicit, now decompose
\(V=V_0\oplus V_1\oplus V_2\) and write
\[
 T=\begin{pmatrix}
 A&B_1&B_2\\C_1&D_{11}&D_{12}\\C_2&D_{21}&D_{22}
 \end{pmatrix},\qquad Q_2=(I-uD_{22})^{-1}.
\]
Eliminating the third component leaves the four blocks
\begin{equation}
 \begin{aligned}
 A'&=A+uB_2Q_2C_2,&
 B_1'&=B_1+uB_2Q_2D_{21},\\
 C_1'&=C_1+uD_{12}Q_2C_2,&
 D_{11}'&=D_{11}+uD_{12}Q_2D_{21}.
 \end{aligned}
 \label{mem:schur-intermedio}
\end{equation}

\begin{proposition}[Transitivity of elimination]
If \(D_*=(D_{ij})_{i,j=1}^2\), the complement obtained in two stages
coincides with that obtained by eliminating both memories together:
\begin{equation}
 \begin{split}
 I-uA'-u^2B_1'(I-uD_{11}')^{-1}C_1'
 ={}&I-uA\\
 &-u^2(B_1\ B_2)(I-uD_*)^{-1}\binom{C_1}{C_2}.
 \end{split}
 \label{mem:schur-transitivo}
\end{equation}
\end{proposition}
\begin{proof}
Solving the third equation of \((I-uT)w=v\) and substituting it into
the first two produces \eqref{mem:schur-intermedio}. Eliminating the
second component then gives the left-hand side of
\eqref{mem:schur-transitivo}. Solving the two memory equations
together gives the right-hand side, by Proposition~\ref{mem:prop-schur}.
The formal solution is unique because \(I-uT\) has constant term
\(I\). Both procedures also transform the source and reader according
to \eqref{mem:schur-fuente}--\eqref{mem:schur-lector}.
\end{proof}

Balance \eqref{mem:balance} describes record increments.
The preceding identities allow them to be transported when memory is
reused in dodecaphase closure; they do not by themselves determine
additional analytic coefficients.

